\section{Quantum Reservoirs and Reset-Window Models}
\label{sec:models}

A QRC carries information from the past in its quantum state.
A reset-window QELM receives recent history explicitly through its input.
We ask when this input window contains enough information to reproduce
the QRC's output.

A QRC processes inputs one at a time. At discrete step $t$, each new input
$a_t$ updates the state left by the previous step, and selected measurements
provide the output \cite{fujii2017harnessing,mujal2021opportunities}:
\[
 \rho_t=\Phi_{a_t}(\rho_{t-1}),
 \qquad y_t=\mathcal C(\rho_t).
\]
Here, $\rho_t$ is the carried quantum state, and the quantum channel
$\Phi_{a_t}$ describes its physical evolution, including dissipation.
The output map $\mathcal C$ collects a fixed finite set of observable
expectations $\operatorname{Tr}(O_j\rho_t)$, $j=1,\ldots,m$.
These are the average values of chosen measurements; $m$ is the output
dimension. Angle brackets in Figure~\ref{fig:information-interface} denote
these expectations. A trained linear readout can also be included in
$\mathcal C$ by absorbing its fixed weights into the observables.

A reset-window QELM accesses the past through the inputs supplied to each
prediction. It starts from the same quantum state each time and encodes a
window of the $d$ most recent inputs
\cite{innocenti2023potential,kawanabe2026tdqelm}.
For a history $h$ ending at time $t$, with target output $y(h)=y_t$,
this window is
\[
 S_d(h)=(a_{t-d+1},\ldots,a_t).
\]
The encoder $E(S_d(h))$ transforms this window into quantum features,
obtained from observable expectations of $P_j$, $j=1,\ldots,B$,
where $j$ also distinguishes measurement checkpoints.
Thus $B$ counts measured expectations, not qubits.
A trained linear readout combines these features. With exact expectations,
the complete prediction can therefore be written as
\[
 \widehat y_t=f(S_d(h)),
\]
where $f$ includes both the quantum feature map and the readout.
The quantum system starts afresh for every prediction, while the input
window supplies the temporal information.

Identical windows therefore produce identical predictions for a
fixed model. Two histories may differ in their earlier inputs, but a
reset-window QELM receives the same information whenever their last $d$
inputs agree. With finitely sampled measurements, individual predictions
can fluctuate, but their probability distributions remain identical.
More qubits or features can enrich the function represented within the
window. They cannot recover an earlier input that the window leaves
unspecified.

Our emulation target is the QRC's chosen output $y_t$, not an external
forecasting target. Matching this output does not require reconstructing
the full quantum state. We hold $\mathcal C$ fixed throughout each
comparison. In our simulations, the outputs are the observable
expectations themselves, so the target readout is the identity.

Figure~\ref{fig:information-interface} illustrates the central question.
The two right-hand plots advance time at a fixed window length. Once an
old input leaves the window, the ideal window prediction cannot use it.
Its effect on the QRC output then fades in one case and persists in the
other. The green prediction is the conditional mean over two equally likely
old-input values with the same later inputs and known baseline.
These are independent illustrations, not evaluation data or fitted QELM
predictions. Decay of one old input with elapsed time does not establish
uniform accuracy as the window length grows; Section~\ref{sec:criterion}
gives that criterion and the exact scalar error bound $D_d/2$. In the input plot,
$k$ indexes displayed steps and $t$ marks the current one.
The next section quantifies how much this missing information limits
emulation.
